Úvodní poznámka: při hodnocení v době dostupné generativní AI se musí řešit současně tři propojené oblasti: transparentnost vůči studentům, ochrana osobních údajů při sběru a zpracování artefaktů a jasná pravidla akademické integrity. Tyto oblasti spolu úzce souvisejí a rozhodnutí v jedné oblasti mají dopad na ostatní; je proto vhodné navrhovat opatření komplexně, nikoli izolovaně. Níže najdete stručné zásady a praktické postupy k okamžitému použití na úrovni předmětu nebo katedry, které lze adaptovat podle konkrétního oboru a typu zadání.

Zásady (konkrétní doporučení)
\begin{itemize}
  \item Požadavek na procesní artefakty: vždy, kde je to pedagogicky oprávněné, vyžadujte od studentů nahrání meziverzí, zápisků postupů, commit logů nebo screenshotů z průběhu práce — procesní důkazy snižují motivaci k neoprávněnému použití AI, protože hodnocení váží postup stejně nebo více než finální výsledek. Praktický příklad: v některých nástrojích lze nastavit povinnost nahrání postupu řešení; u „Pokroku v matematice“ je explicitně možné požadovat nahrání postupu jako součást zadání \cite{kb_ai_101_tipu_pdf_04e4a115_page_8_1}. Procesní artefakty také pomáhají učiteli při ověřování porozumění a identifikaci zdrojů chyb.
  \item Transparentnost vůči studentům: do sylabu uveďte, jaké AI nástroje jsou v kurzu povoleny, za jakých podmínek (povinná dokumentace, omezení generování hotových řešení) a jak mají studenti deklarovat použití AI (general background knowledge). Zveřejněná pravidla by měla být konkrétní a srozumitelná, aby studenti věděli, čeho se mají vyvarovat, a jak dokumentovat oprávněné použití nástrojů.
  \item Minimální sběr osobních údajů: sbírejte jen ta data, která jsou nezbytná pro hodnocení (general background knowledge). Pokud používáte nástroje třetích stran, ověřte smluvní podmínky zpracování (DPA) a možnosti data residency (general background knowledge). Zaznamenávejte, proč jsou jednotlivá data potřeba, a jak dlouho budou uchovávána.
  \item Upřednostněte enterprise/školní licence tam, kde je to možné — integrace do školních balíků (např. Microsoft 365) často nabízí lepší kontrolu nad daty a nastavením soukromí; některé školní licence zahrnují i specifické vzdělávací aplikace jako Minecraft Education nebo funkce OneNote pro převod rukopisu \cite{kb_ai_101_tipu_pdf_04e4a115_page_15_2,kb_ai_101_tipu_pdf_04e4a115_page_14_1}. Tyto licence navíc obvykle umožňují centrální správu přístupů, auditování a případnou integraci do školních systémů LMS.
  \item Jasná pravidla akademické integrity: rozhodněte, které typy úloh AI povolíte (např. povoleno s povinnou reflexí a deklarací) a které jsou zakázané; pro podezření z neférového použití preferujte procesní ověření před automatickými detektory (general background knowledge). Doporučuje se také definovat důsledky porušení pravidel a postupy pro jejich vynucení, včetně postupného stupňování sankcí.
\end{itemize}

Krátký operační checklist pro vyučující (krok za krokem)
\begin{enumerate}
  \item Aktualizujte sylabus: vložte krátký odstavec o pravidlech používání AI, požadavku na deklaraci a popisu požadovaných procesních artefaktů (general background knowledge). Doporučujeme ukázkový text vložit přímo do elektronického sylabu i do první přednášky.
  \item Informujte studenty při zahájení kurzu: vysvětlete prakticky, jak deklarovat použití AI, jaké artefakty nahrávat a proč jsou tyto kroky důležité pro hodnocení a důvěryhodnost. Pořaďte krátký demonstrační návod nebo FAQ v LMS.
  \item Před odevzdáním nastavte požadavek na postup: v LMS vyžadujte nahrání meziverzí, logu nebo exportu z nástrojů (commit history, OneNote stránky apod.) — viz příklady praxe v ostatních vzdělávacích nástrojích, kde je nahrání postupu součástí zadání \cite{kb_ai_101_tipu_pdf_04e4a115_page_8_1}. Explicitní nastavení v LMS usnadní konzistenci hodnocení mezi asistenty a učiteli.
  \item Použijte školní/enterprise řešení pro generování či asistenci (pokud je dostupné), abyste omezili nechtěné sdílení osobních dat do veřejných cloudů (general background knowledge). V rámci Microsoft‑stacku lze využít integrace do Poznámkového bloku / Copilot funkcí pro bezpečnější workflow \cite{kb_ai_101_tipu_pdf_04e4a115_page_14_1}. Ověřte také, zda použitá řešení podporují export auditních záznamů.
  \item Při podezření postupujte systematicky: zkontrolujte procesní artefakty → pohovor se studentem (ověření porozumění) → případná disciplinární opatření podle místních pravidel (general background knowledge). Dokumentujte průběh šetření a uchovejte záznamy pro případné další kroky.
  \item Zpětná vazba a iterace: na konci semestru vyhodnoťte, jak opatření fungovala, sesbírejte názory studentů a asistentů a upravte postupy pro další běh kurzu. Takto udržíte pravidla aktuální vůči rychle se měnícím technologiím.
\end{enumerate}

Vzorový text do sylabu (kratká formulace — upravte dle potřeby)
\begin{quote}
Studenti mohou v této části kurzu využívat generativní AI pouze za podmínky, že: (1) každé použití AI bude v odevzdání explicitně deklarováno, (2) k odevzdání budou přiloženy procesní artefakty (meziverze, log změn, krátké vysvětlení postupu) a (3) konečné hodnocení zahrnuje ověření osobního porozumění učitelem. Neoprávněné použití AI bude posuzováno podle pravidel akademické integrity fakulty. V případě nejasností kontaktujte vyučujícího nebo garanta předmětu. (general background knowledge)
\end{quote}

Krátké technické doporučení (bezpečnost dat a nástroje)
\begin{itemize}
  \item Pokud ve výuce využíváte nástroje, které umí převádět rukopis nebo pomáhat s řešením (např. OneNote či specializované „education“ aplikace), nastavte to tak, aby citlivá data studentů nesdílela více než nutné a aby bylo jasné, kde jsou data uchovávána \cite{kb_ai_101_tipu_pdf_04e4a115_page_8_1,kb_ai_101_tipu_pdf_04e4a115_page_14_1}. Zvažte anonymizaci vzorků tam, kde to hodnocení umožňuje.
  \item U didaktických světů a simulací (např. Minecraft Education) zkontrolujte licenční podmínky a dostupnost bezplatných hodin kódu jako pomůcky pro výuku AI principů; tyto prostředí mohou být užitečné pro demonstraci etických aspektů a principů práce s daty \cite{kb_ai_101_tipu_pdf_04e4a115_page_15_2}. Pro simulace také ověřte, zda prostředí neposkytuje zbytečně široké sdílení studentových údajů.
  \item Správa přístupů a šablony: vytvořte jednoduché šablony pro nahrávání procesních artefaktů a instrukce pro pojmenování souborů, aby bylo možné rychle vyhledávat a kontrolovat odevzdání v LMS.
\end{itemize}

Poznámka o právních a GDPR aspektech: konkrétní povinnosti (např. provedení DPIA, dohody o zpracování osobních údajů, role správce/zpracovatele) závisí na rozsahu zpracování dat a na lokálních pravidlech VUT — tyto kroky považujeme za nezbytné, ale jejich podrobný postup je uveden v kapitole o ochraně osobních údajů a v přiložené šabloně DPIA (viz Přílohy). (general background knowledge) V praxi je vhodné konzultovat právní oddělení instituce před zavedením nových nástrojů do výuky a uchovávat záznamy o provedených posouzeních.

Shrnutí: kombinujte pedagogicky odůvodněné požadavky na procesní artefakty s jasnou komunikací pravidel do sylabu a s preferencí školních/enterprise řešení tam, kde to zajišťuje lepší kontrolu nad daty. Tam, kde to technicky umožňují nástroje, využijte možnosti nastavení povinného nahrání postupu řešení jako standardního požadavku při hodnocení \cite{kb_ai_101_tipu_pdf_04e4a115_page_8_1}. Pravidelně revidujte pravidla a technická nastavení podle zkušeností z praxe a změn v legislativě či nabídce nástrojů.